\par\noindent\textbf{(1)} 概率流 \(J=k_{01}p_0-k_{10}p_1\)，故 \(\dot p_1=J=- (k_{01}+k_{10})p_1+k_{01}\)。平衡概率 \(\pi_1=k_{01}/(k_{01}+k_{10})\)，\(\pi_0=k_{10}/(k_{01}+k_{10})\)，在单一边连接的二态网络必满足稳态详细平衡。
\par\noindent\textbf{(2)} 线性常微分方程解为 \(p_1(t)=\pi_1+[p_1(0)-\pi_1]e^{-(k_{01}+k_{10})t}\)，故弛豫时间 \(\tau=1/(k_{01}+k_{10})\)。这与两个逆向单独等待时间不同，不能将 \(\tau\) 误写为 \(1/k_{01}\)。
\par\noindent\textbf{(3)} 定义 \(D(p\Vert\pi)=\sum_ip_i\ln(p_i/\pi_i)\)。对主方程求导得 \(-dD/dt=J\ln[(k_{01}p_0)/(k_{10}p_1)]\)。利用 \((x-y)\ln(x/y)\ge0\) 得 \(\sigma/k_B=-\dot D\ge0\)，等号仅在净概率流为零的平衡态成立。
\par\noindent\textbf{(4)} 在仅两态且仅一对逆向通道的稳态，\(\dot p_1=J=0\) 强制净流为零。若要在定常概率分布中保持非零循环流，必须有至少一个闭合回路，比如三态环且正反路径率乘积不等；否则不能声称稳态双态系统持续耗散正熵。
